% -------------------------------------------------------------------------
% Block-size sweep write-up (ΔLVR/ΔFees)
% Based on: docs/LVR_explanation.md and the outputs produced by LVR_vs_blocksize.py
% Location of artifacts (this run):
%   images/
%
% This text is meant to be pasted into a LaTeX report (requires: amsmath, amssymb, graphicx, booktabs).
% -------------------------------------------------------------------------
\documentclass[11pt]{article}

\usepackage[margin=1in]{geometry}
\usepackage{amsmath,amssymb,amsfonts,amsthm}
\usepackage{mathtools}
\usepackage{bm}
\usepackage{hyperref}
\usepackage{graphicx}
\usepackage{enumitem}
\usepackage{booktabs}
\usepackage{microtype}
\usepackage{tcolorbox}
\usepackage{tabularx}
\usepackage{caption}
\usepackage[table]{xcolor}
\definecolor{groupgray}{RGB}{235,235,235} % change to any solid color you like


% Define a custom environment that behaves like a beamer block
\newtcolorbox{block}[1]{
  colback=blue!5!white,
  colframe=blue!75!black,
  title=#1
}
\usepackage{authblk}
\title{Agent-Based Modelling of Uniswap v3:\\
Microstructure, Dynamic Fees, and Loss-Versus-Rebalancing}
\author[1]{Daniele Maria Di Nosse\textsuperscript{\dag}}
\author[1]{Fabrizio Lillo}

\affil[1]{Scuola Normale Superiore, Pisa, Italy}
\date{\today}

\begin{document}

\section{Block-size sweep: the fee-coverage ratio \texorpdfstring{$\Delta\mathrm{LVR}/\Delta\mathrm{Fees}$}{ΔLVR/ΔFees}}

\subsection{Metric definitions and accounting identity}

Loss-versus-Rebalancing (LVR) is a structural cost of AMM liquidity provision: when the external reference price
moves, the AMM quote becomes stale and arbitrage trades realign the pool to the external market, extracting value
from liquidity providers. Following the accounting used in this repository (see \texttt{docs/LVR\_explanation.md}),
LP wealth decomposes as
\begin{equation}
  V^{LP}_t \;=\; V^{reb}_t \;+\; F_t \;-\; \mathrm{LVR}_t,
\end{equation}
where:
\begin{itemize}
  \item $V^{LP}_t$ is the LP's mark-to-market wealth valued at the end-of-block CEX mid price $m_t$;
  \item $V^{reb}_t$ is the \emph{rebalancing benchmark} value: it holds the same inventory path as the LP but executes
        inventory adjustments at the external price $m_t$ (frictionless execution);
  \item $F_t$ is cumulative fee value (in token1 units, valued at $m_t$);
  \item $\mathrm{LVR}_t$ is cumulative loss-versus-rebalancing.
\end{itemize}

Define the \emph{hedged PnL} as
\begin{equation}
  \mathrm{PnL}^{hedged}_t \;:=\; V^{LP}_t - V^{reb}_t.
\end{equation}
Then the identity above implies
\begin{equation}
  \mathrm{PnL}^{hedged}_t \;=\; F_t - \mathrm{LVR}_t.
  \label{eq:hedged_pnl_identity}
\end{equation}
This is the key interpretation tool: on a hedged basis, fees are the ``revenue'' and LVR is the ``structural cost''.

\subsection{Per-block increments and the ratio plotted in the sweep}

The script \texttt{LVR\_vs\_blocksize.py} converts cumulative series into \emph{per-block increments} after discarding a
burn-in prefix of \texttt{skip\_step} blocks. For any cumulative series $X_t$, define
\begin{equation}
  \Delta X_t := X_t - X_{t-1}.
\end{equation}
Using \eqref{eq:hedged_pnl_identity}, per block we have
\begin{equation}
  \Delta\mathrm{PnL}^{hedged}_t \;=\; \Delta\mathrm{Fees}_t \;-\; \Delta\mathrm{LVR}_t.
\end{equation}

In this experiment, the denominator is the \emph{fee-flow value in the block}, computed as
\begin{equation}
  \Delta\mathrm{Fees}_t
  \;:=\;
  (\Delta\mathrm{fees0}_t)\,m_t + (\Delta\mathrm{fees1}_t),
\end{equation}
i.e.\ in-block token0 fee flow is valued at the \emph{end-of-block} CEX mid price $m_t$ and added to the in-block token1
fee flow. This is the script's \texttt{fee\_definition=flow} mode and intentionally avoids any mark-to-market revaluation
of \emph{previously accumulated} token0 fees.

For each agent cohort, the plotted ``fee coverage'' ratio is
\begin{equation}
  r_t \;:=\; \frac{\Delta\mathrm{LVR}_t}{\Delta\mathrm{Fees}_t}.
  \label{eq:ratio_def}
\end{equation}
Combining with the per-block identity gives the most direct interpretation:
\begin{equation}
  \Delta\mathrm{PnL}^{hedged}_t \;=\; (1-r_t)\,\Delta\mathrm{Fees}_t.
  \label{eq:ratio_to_hedged_pnl}
\end{equation}

\paragraph{Interpretation of the ratio.}
Equation \eqref{eq:ratio_to_hedged_pnl} turns the ratio into a simple economic diagnostic:
\begin{itemize}
  \item If $r_t>1$, then $(1-r_t)<0$ and the block-level hedged PnL is negative:
        \[
          \Delta\mathrm{PnL}^{hedged}_t < 0.
        \]
        Fees \emph{do not cover} the structural adverse-selection cost (LVR). Liquidity provision is unprofitable
        on a hedged basis in such blocks.
  \item If $0<r_t<1$, then the block-level hedged PnL is positive:
        \[
          \Delta\mathrm{PnL}^{hedged}_t > 0.
        \]
        Fees cover LVR, and $(1-r_t)$ is the fraction of fees that remains as hedged profit in that block.
        For example, $r_t=0.8$ means $\sim 20\%$ of fee revenue remains after paying the structural LVR cost.
  \item If $r_t=1$, fees exactly offset LVR and hedged PnL is approximately zero (break-even on a hedged basis).
  \item If $r_t<0$, then $\Delta\mathrm{LVR}_t<0$ (given $\Delta\mathrm{Fees}_t>0$ in the ratio computation),
        which implies
        \[
          \Delta\mathrm{PnL}^{hedged}_t > \Delta\mathrm{Fees}_t.
        \]
        This is a \emph{negative-LVR} regime: the agent's realized execution path (relative to the rebalancing benchmark)
        is \emph{better than} frictionless trading at $m_t$ for the same inventory changes. For passive LPs this is unusual,
        but for \emph{active or selectively-timed} strategies (e.g.\ JIT liquidity) it can occur because the strategy chooses
        when to be exposed and can capture favorable selection.
\end{itemize}

\paragraph{Important sampling note (denominator filtering).}
The script masks blocks where $\Delta\mathrm{Fees}_t$ is non-positive or numerically tiny (to avoid unstable ratios).
Therefore, the ratio distribution is conditional on ``blocks with positive fee flow.'' This is typically appropriate because
blocks with no trades (and hence no fees) do not meaningfully define a fee-coverage ratio.

\subsection{Cohorts and Jiter-specific adjustment}

The sweep produces results for three cohorts:
\begin{itemize}
  \item \textbf{Passive LPs:} baseline LPs with wide, relatively static positions.
  \item \textbf{Active LPs:} programmatic/repositioning LPs with narrower liquidity and more frequent recentering.
  \item \textbf{Jiter (JIT LP):} a liquidity provider that mints and burns around large swaps within a block.
\end{itemize}

For the Jiter, the script adjusts LVR to account for explicit flash-loan financing costs (so we do not misattribute financing
expense to ``structural LVR''):
\begin{equation}
  \Delta\mathrm{LVR}^{jit}_t
  \;:=\;
  \Delta\mathrm{Fees}^{jit}_t
  \;-\;
  \left(\Delta\mathrm{PnL}^{hedged,jit}_t + \Delta\mathrm{FlashPaid}_t\right).
  \label{eq:jiter_lvr_def}
\end{equation}
The corresponding ratio is $r^{jit}_t=\Delta\mathrm{LVR}^{jit}_t/\Delta\mathrm{Fees}^{jit}_t$.

\paragraph{Selection for ``successful JIT'' blocks.}
When plotting Jiter distributions, the script additionally filters to blocks in which a JIT cycle successfully surrounds an
executed swap (recovered from \texttt{jiter\_activity\_cum}). This is a deliberate choice: it focuses the analysis on the
economics of \emph{executed JIT opportunities} rather than averaging over idle blocks where the strategy did nothing.
However, this also introduces a potential \emph{selection bias} if successful JIT events are rare (see discussion in the
static-fee results below).

\subsection{Experiment structure and outputs}

The independent variable is the block size $B=\texttt{block\_time}$, swept from $B=2$ to $B=16$ (inclusive). For each $B$,
the script runs multiple simulations with different random seeds, converts cumulative series into $\Delta$-series after burn-in,
and pools the samples across time and seeds.

Artifacts written by the sweep include:
\begin{itemize}
  \item \textbf{NPZ bundles} (\texttt{dLVR\_arrays\_*.npz}): raw 3D arrays indexed by $(B,\text{seed},t)$ for both $\Delta\mathrm{LVR}$
        and $\Delta\mathrm{LVR}/\Delta\mathrm{Fees}$;
  \item \textbf{CSV summaries} (\texttt{dLVR\_over\_dFees\_summary\_*.csv}): pooled distribution statistics per cohort and $B$
        (median, percentiles, etc.);
  \item \textbf{Plots} in \textbf{PNG} and interactive \textbf{HTML} (Plotly): the median ratio vs $B$ with error bars showing the
        central $95\%$ interval of the pooled samples.
\end{itemize}

% -------------------------------------------------------------------------
\section{Results: static fee mode (fixed fee)}

In the \textbf{static} fee mode, the protocol fee is constant (no feedback from toxicity). The key empirical feature in these
plots is that the LP fee coverage ratio increases strongly with block size: as blocks become ``longer,'' LVR per block grows
faster than fee revenue per block.

\subsection{Only passive LPs}

% (Replace filename if needed; PIDs differ across runs.)
\begin{figure}[h]
  \centering
  \includegraphics[width=\linewidth]{images/dLVR_over_dFees_medians_vs_block_time_static_flow_test_runs50_pid42272.png}
  \caption{Static fees, passive LP only: median and 95\% interval of $r=\Delta\mathrm{LVR}/\Delta\mathrm{Fees}$ vs block size $B$.}
  \label{fig:static_passive_only_ratio}
\end{figure}

Passive LPs exhibit $r>1$ across all $B$, and the median ratio grows with $B$.
In the CSV summary (\texttt{.../dLVR\_over\_dFees\_summary\_static\_flow\_test\_runs50\_pid42272.csv}),
the median is approximately
\[
  r_{\mathrm{passive}} \approx 1.26 \text{ at } B=2,
  \qquad
  r_{\mathrm{passive}} \approx 3.50 \text{ at } B=16.
\]
From \eqref{eq:ratio_to_hedged_pnl}, the typical block has negative hedged PnL, and this worsens with $B$:
\[
  \Delta\mathrm{PnL}^{hedged}_t
  \approx
  (1-1.26)\Delta\mathrm{Fees}_t \approx -0.26\,\Delta\mathrm{Fees}_t \quad (B=2),
\]
\[
  \Delta\mathrm{PnL}^{hedged}_t
  \approx
  (1-3.50)\Delta\mathrm{Fees}_t \approx -2.50\,\Delta\mathrm{Fees}_t \quad (B=16).
\]
Economically, passive LPs are systematically exposed to adverse selection: when the reference price moves within a block, the AMM
quote becomes stale for longer, arbitrage extracts more value, and fixed fees are insufficient to compensate.

\subsection{Passive + active LPs}

\begin{figure}[h]
  \centering
  \includegraphics[width=\linewidth]{images/dLVR_over_dFees_medians_vs_block_time_static_flow_test_runs50_pid42482.png}
  \caption{Static fees, active and passive LPs: median and 95\% interval of $r=\Delta\mathrm{LVR}/\Delta\mathrm{Fees}$ vs $B$.}
  \label{fig:static_active_passive_ratio}
\end{figure}

Both cohorts again have $r>1$, and both rise with $B$. Quantitatively (from
\texttt{.../dLVR\_over\_dFees\_summary\_static\_flow\_test\_runs50\_pid42482.csv}):
\[
  r_{\mathrm{active}} \approx 1.57 \text{ at } B=2,\quad r_{\mathrm{active}} \approx 4.14 \text{ at } B=16,
\]
\[
  r_{\mathrm{passive}} \approx 1.53 \text{ at } B=2,\quad r_{\mathrm{passive}} \approx 3.96 \text{ at } B=16.
\]

The economic message is that \textbf{concentration increases LVR exposure per unit fee} under a fixed fee schedule.
Active LPs (narrower liquidity, more concentrated marginal liquidity) are traded against more aggressively when the price moves,
which increases the ``pick-off'' component of returns. Without a compensating fee increase, the active LP cohort exhibits the
largest fee coverage ratios (and hence the worst hedged profitability).

\subsection{Passive + active LPs + Jiter}

\begin{figure}[h]
  \centering
  \includegraphics[width=\linewidth]{images/dLVR_over_dFees_medians_vs_block_time_static_flow_test_runs50_pid42657.png}
  \caption{Static fees, active+passive LPs and Jiter: median and 95\% interval of $r=\Delta\mathrm{LVR}/\Delta\mathrm{Fees}$ vs $B$.}
  \label{fig:static_with_jiter_ratio}
\end{figure}

The LP cohorts remain in the $r>1$ regime and continue to rise with $B$ (consistent with the previous figure).
The new feature is the \textbf{Jiter's negative ratio}, with median roughly
\[
  r_{\mathrm{jit}} \approx -5.8 \text{ at } B=2,
  \qquad
  r_{\mathrm{jit}} \approx -8.8 \text{ at } B=16
\]
(from \texttt{.../dLVR\_over\_dFees\_summary\_static\_flow\_test\_runs50\_pid42657.csv}).

By \eqref{eq:ratio_to_hedged_pnl} and the Jiter definition \eqref{eq:jiter_lvr_def}, $r_{\mathrm{jit}}<0$ implies that in these
blocks the Jiter's hedged PnL (net of flash costs) exceeds fee revenue:
\[
  \Delta\mathrm{PnL}^{hedged,jit}_t + \Delta\mathrm{FlashPaid}_t
  \;>\;
  \Delta\mathrm{Fees}^{jit}_t,
\]
equivalently $\Delta\mathrm{LVR}^{jit}_t<0$.

\paragraph{Economic interpretation.}
Negative ``LVR'' should be read as \textbf{favorable selection} for an \emph{actively timed} strategy: conditional on entering only
around specific swaps, the Jiter can earn fees while avoiding most inventory risk and avoiding being present during adverse
arbitrage episodes. In that sense, the strategy can outperform the frictionless rebalancing benchmark given the same realized
inventory changes.

\paragraph{Selection bias warning (crucial in the static-fee case).}
In the static-fee run, the Jiter statistics are computed only on \emph{successful JIT blocks} (where a mint/burn roundtrip surrounds
an executed swap), and the number of such blocks is extremely small (e.g.\ $n\approx 5$--$17$ pooled samples per $B$ in the CSV).
This induces a \textbf{strong selection bias}: the sample is dominated by a handful of rare, extremely profitable JIT opportunities.
If the strategy attempts JIT frequently but only succeeds when conditions are especially favorable, then conditioning on success
selects precisely those extreme outcomes, which can easily push the measured $\Delta\mathrm{LVR}^{jit}$ negative and make
$r_{\mathrm{jit}}$ strongly negative.
In other words, in the static-fee configuration the Jiter plot is best interpreted as:
\begin{quote}
  \emph{``When a JIT opportunity actually executes, it tends to be very profitable relative to the benchmark''}
\end{quote}
rather than as an unconditional, time-average statement about all blocks.

% -------------------------------------------------------------------------
\section{Results: toxicity fee mode (fee responds to toxicity)}

In the \textbf{toxicity} fee mode, the fee responds to a toxicity/basis signal. The key empirical feature is that the fee coverage
ratio becomes \textbf{stable and below one} for LPs: fees adapt upward in precisely those environments where LVR pressure is high,
restoring hedged profitability.

\subsection{Only passive LPs}

\begin{figure}[h]
  \centering
  \includegraphics[width=\linewidth]{images/dLVR_over_dFees_medians_vs_block_time_toxicity_flow_test_runs50_pid13640.png}
  \caption{Toxicity fees, passive LP only: median and 95\% interval of $r=\Delta\mathrm{LVR}/\Delta\mathrm{Fees}$ vs $B$.}
  \label{fig:toxicity_passive_only_ratio}
\end{figure}

Passive LPs exhibit $r<1$ and nearly flat in $B$. From
\texttt{.../dLVR\_over\_dFees\_summary\_toxicity\_flow\_test\_runs50\_pid13640.csv}, the median is approximately
\[
  r_{\mathrm{passive}} \approx 0.74 \text{ for } B=2\ldots 16.
\]
Thus the typical block has positive hedged profitability:
\[
  \Delta\mathrm{PnL}^{hedged}_t \approx (1-0.74)\Delta\mathrm{Fees}_t \approx 0.26\,\Delta\mathrm{Fees}_t.
\]
Economically, the toxicity controller is doing what it is designed to do: increase fee revenue in environments where adverse
selection (and hence LVR) is elevated, so that LPs are compensated on a hedged basis.

\subsection{Passive + active LPs}

\begin{figure}[h]
  \centering
  \includegraphics[width=\linewidth]{images/dLVR_over_dFees_medians_vs_block_time_toxicity_flow_test_runs50_pid13959.png}
  \caption{Toxicity fees, active and passive LPs: median and 95\% interval of $r=\Delta\mathrm{LVR}/\Delta\mathrm{Fees}$ vs $B$.}
  \label{fig:toxicity_active_passive_ratio}
\end{figure}

Both cohorts remain below one, with medians roughly in the $0.8$--$0.9$ range (see
\texttt{.../dLVR\_over\_dFees\_summary\_toxicity\_flow\_test\_runs50\_pid13959.csv}):
\[
  r_{\mathrm{active}} \approx 0.87 \text{ at } B=2,\quad r_{\mathrm{active}} \approx 0.83 \text{ at } B=16,
\]
\[
  r_{\mathrm{passive}} \approx 0.85 \text{ at } B=2,\quad r_{\mathrm{passive}} \approx 0.81 \text{ at } B=16.
\]
The economic message is that fee adaptivity largely offsets the greater adverse selection exposure of concentrated liquidity.
Active LPs still tend to have slightly larger $r$ (more LVR per unit fee), but the controller keeps them in a regime where fees cover
LVR in the median block.

\subsection{Passive + active LPs + Jiter}

\begin{figure}[h]
  \centering
  \includegraphics[width=\linewidth]{images/dLVR_over_dFees_medians_vs_block_time_toxicity_flow_test_runs50_pid14207.png}
  \caption{Toxicity fees, active+passive LPs and Jiter: median and 95\% interval of $r=\Delta\mathrm{LVR}/\Delta\mathrm{Fees}$ vs $B$.}
  \label{fig:toxicity_with_jiter_ratio}
\end{figure}

LP cohorts remain near $r\approx 0.84$--$0.87$.
The Jiter shows a lower median ratio, around
\[
  r_{\mathrm{jit}} \approx 0.57 \text{ at } B=2,\quad r_{\mathrm{jit}} \approx 0.54 \text{ at } B=16
\]
(from \texttt{.../dLVR\_over\_dFees\_summary\_toxicity\_flow\_test\_runs50\_pid14207.csv}).

Economically, $r_{\mathrm{jit}}<1$ indicates that JIT liquidity keeps a larger fraction of fees as hedged profit (net of flash costs)
than the average LP cohort:
\[
  \Delta\mathrm{PnL}^{hedged,jit}_t
  \approx
  (1-r_{\mathrm{jit}})\Delta\mathrm{Fees}^{jit}_t
  \;\approx\;
  0.43\,\Delta\mathrm{Fees}^{jit}_t
  \quad (\text{median}).
\]
Unlike the static-fee case, the Jiter ratio here is positive (no systematic negative-LVR regime), and the sample size of successful
JIT events is large, so selection effects are much weaker and the estimate is more representative of typical JIT performance when it
does execute.

% -------------------------------------------------------------------------
\section{Why does the ratio grow roughly linearly with block size?}

Two complementary explanations are useful: one theoretical (diffusion scaling), and one model-specific (how activity is parameterized).

\subsection{Diffusion intuition: LVR grows with elapsed time between quote updates}

In the LVR theory, when the external reference price follows a diffusion with instantaneous volatility $\sigma_t$, LVR accrues at a
rate proportional to $\sigma_t^2$ and marginal liquidity. In particular, in continuous time LVR over an interval is an integral of a
volatility-squared term (see \texttt{docs/LVR\_explanation.md} for the Uniswap v3 specialization).

Holding other factors fixed, the important scaling is:
\begin{equation}
  \mathbb{E}[\mathrm{LVR} \text{ over an interval of length } \Delta t]
  \;\propto\;
  \sigma^2 \,\Delta t.
  \label{eq:lvr_time_scaling}
\end{equation}
In the simulator, increasing $B=\texttt{block\_time}$ makes each block represent a longer interval before the ``end-of-block'' price
$m_t$ is realized and before the pool is compared against it. Larger blocks therefore produce larger typical price changes within
a block (variance scales with time), which increases the magnitude of mispricing and arbitrage extraction in that block, consistent
with \eqref{eq:lvr_time_scaling}. This naturally produces approximately linear growth of $\Delta\mathrm{LVR}$ with $B$.

\subsection{Model-specific: fees per block need not scale with block duration}

The ratio is $\Delta\mathrm{LVR}/\Delta\mathrm{Fees}$. Even if $\Delta\mathrm{LVR}$ grows with $B$, the ratio will only remain stable if
fee revenue also grows proportionally.

In this ABM configuration, trader arrivals are specified largely \emph{per block} (e.g.\ expected number of trade intents per block),
not per unit of continuous time. As a result, increasing $B$ can raise LVR per block (because each block contains a longer exposure
to price diffusion and stale quotes) without mechanically increasing the number of fee-generating trades per block by the same factor.
Under \textbf{static fees}, the fee rate does not respond, so $\Delta\mathrm{Fees}$ grows weakly compared to $\Delta\mathrm{LVR}$ and the
ratio rises with $B$.

Under \textbf{toxicity fees}, by construction the protocol fee increases when the environment is more toxic (larger basis / more
adverse selection). Increasing $B$ tends to increase measured toxicity because quotes are stale for longer and deviations between DEX
and reference price become larger. The controller therefore increases $\Delta\mathrm{Fees}$ in the same regimes where $\Delta\mathrm{LVR}$
would otherwise rise, keeping the ratio comparatively flat and below one.

\subsection{Heavy tails and why medians matter}

Because the ratio divides by $\Delta\mathrm{Fees}$, blocks with very small (but positive) fee flows can generate extremely large ratios
even if $\Delta\mathrm{LVR}$ is moderate. This creates heavy-tailed distributions and makes means unstable.
The plots therefore emphasize medians (and 95\% intervals) as robust summaries of ``typical'' fee coverage.

% -------------------------------------------------------------------------
\section{Summary of economic takeaways}

\begin{itemize}
  \item \textbf{Static fee mode:} both passive and active LP cohorts have $\Delta\mathrm{LVR}/\Delta\mathrm{Fees}>1$, increasingly so at
        larger $B$. The structural adverse selection cost dominates fixed fee revenue, producing negative hedged profitability.
        Active (more concentrated) LPs tend to have the worst coverage ratios.
  \item \textbf{Toxicity fee mode:} the ratio drops below one and becomes much more stable across $B$. This indicates the controller
        is broadly successful at matching fee revenue to adverse selection intensity, restoring positive hedged profitability for LPs.
  \item \textbf{Jiter:} in toxicity mode the Jiter has $0<\Delta\mathrm{LVR}/\Delta\mathrm{Fees}<1$ and typically keeps a large fraction
        of fees as hedged profit (net of flash costs). In static mode the Jiter ratio appears strongly negative, but this is best
        understood as a consequence of (i) selective timing, and (ii) strong selection bias from conditioning on a tiny number of
        successful JIT executions: only a handful of extremely profitable JIT opportunities appear in the sample, pushing measured
        $\Delta\mathrm{LVR}$ negative.
\end{itemize}

\end{document}